\begin{claim}{D.74}{Proved}{Full section centralizers from twisted relative duality}
Retain the actual normalized involution, full profinite group,
and labelled extension of D.34:
\[
 \begin{gathered}
 F=\widehat F(x,y),\qquad z=(xy)^{-1},\qquad
 E=F\rtimes_\alpha\langle s\rangle,\\
 \alpha(x)=x^{-1},\qquad \alpha(y)=f^{-1}y^{-1}f,
 \qquad f\in F'.
 \end{gathered}
 \tag{D.74.1}
\]
Thus $\alpha^2=1$, and its multiplier at every prime is
the global element $-1$.  Put
\[
 s_0=s,\qquad s_1=x^{e_2}s,\qquad s_2=y^{e_2}fs,
 \qquad S_i=\langle s_i\rangle,\qquad C_i=C_F(s_i).
 \tag{D.74.2}
\]
Then
\[
                 C_i=1,\qquad N_E(S_i)=S_i
                       \quad(0\leq i<3).
 \tag{D.74.3}
\]
These equalities concern the full profinite groups.  The
extension $E$ is real projective in Haran's sense, with the
actual finite embedding-problem formulation given below.
There is exactly one $h\in F$ satisfying
$f=h^{-1}\alpha(h)$; it lies in $F'$.
For $\alpha\in GT_0$, this same unique element satisfies
$h\theta(h)=1$, where $\theta$ exchanges $x$ and $y$.

The proof gives an exact finite-coefficient boundary
restriction calculation.  D.35's remaining peripheral
conjugator problem is the displayed equation (D.74.40)
with this unique $h$.  Its solution and a free-product
presentation of $E$ are not assertions of this result.
\end{claim}
\begin{proof}
\emph{The three actual peripheral extensions.}
Let $L_1=\overline{\langle x\rangle}$,
$L_2=\overline{\langle y\rangle}$, and
$L_3=\overline{\langle z\rangle}$, and set
$D_j=N_E(L_j)$.  A.7 and D.34 give the exact maps
\[
 \begin{split}
 1\longrightarrow L_j\longrightarrow D_j
       &\xrightarrow{\ u\ }C_2\longrightarrow1,\\
 \widehat{\mathbb Z}\rtimes_{-1}C_2
       &\longrightarrow D_j,\qquad
       (a,\tau^\epsilon)\longmapsto c_j^a b_j^\epsilon ,
 \end{split}
 \tag{D.74.4}
\]
where $(c_1,c_2,c_3)=(x,y,z)$ and $b_j$ is an actual
involution in $D_j$ with $u(b_j)=1$.
The inverse of the second map reads $\epsilon=u(d)$ and
then the unique exponent of $d b_j^{-\epsilon}\in L_j$.
Here and throughout, $u:E\to C_2=\mathbb F_2$ is the
specified quotient map.

The two $D_j$-conjugacy classes of subgroups of order two
are represented by $\langle b_j\rangle$ and
$\langle c_j^{e_2}b_j\rangle$.  Each has normalizer
equal to itself in $D_j$: an element of $L_j$ commuting
with either involution has exponent $a$ with $2a=0$,
and hence is the identity.
Every finite subgroup of $E$ injects into $E/F$, because
$F$ is torsion-free.  D.34 identifies all three
$E$-conjugacy classes of nontrivial finite subgroups as
$S_0,S_1,S_2$ and proves that the six occurrences just
listed form the endpoints of a triangle.  Consequently
each $S_i$ occurs exactly twice among those six classes.

The full normalizers are already specified without any
assumption on $C_i$.  There are inverse continuous maps
\[
 \begin{aligned}
 C_i\times S_i&\longrightarrow N_E(S_i),
                         &(h,s_i^\epsilon)&\longmapsto hs_i^\epsilon,\\
 N_E(S_i)&\longrightarrow C_i\times S_i,
                         &w&\longmapsto(ws_i^{-\epsilon},s_i^\epsilon),
                         \quad \epsilon=u(w).
 \end{aligned}
 \tag{D.74.5}
\]
A normalizer fixes the unique nonidentity element of
$S_i$, which proves that the second formula has its
stated target.  These are the normalizer maps that will
occur in every restriction below.

\emph{One actual coefficient module for the relative theory.}
Write $k=\mathbb F_2$.  Form the compact permutation
$k[[E]]$-module and its augmentation kernel
\[
 P=\bigoplus_{j=1}^3 k[[E/D_j]],\qquad
 0\longrightarrow\Delta\longrightarrow P
             \xrightarrow{\ \varepsilon\ } k\longrightarrow0 .
 \tag{D.74.6}
\]
There is an actual $F$-equivariant homeomorphism
\[
             F/L_j\longrightarrow E/D_j,\qquad
                        aL_j\longmapsto aD_j .
 \tag{D.74.7}
\]
It is injective since $F\cap D_j=L_j$ and surjective
since $E=FD_j$; compactness gives its inverse continuity.
Thus $\Delta|_F$ is precisely the augmentation kernel
for the three specified full peripheral subgroups of
$F$.  The additional $E$-action is retained, rather
than replacing it by an action fixing the $L_j$
pointwise.  In the coordinates of (D.74.7), it is
\[
       e\cdot(aL_j)=e a b_j^{-u(e)}L_j .
 \tag{D.74.8}
\]
The displayed representative lies in $F$.  Changing
$a$ by $a\ell$, $\ell\in L_j$, gives the same coset
because $b_j$ normalizes $L_j$.  This is the action
transported from $E/D_j$, so it obeys the group law.

For a finite continuous $E$-module $M$ over $k$, put
\[
 \mathscr H_M=\operatorname{Hom}_{k,\mathrm{cont}}(\Delta,M),
 \qquad (e\varphi)(d)=e\varphi(e^{-1}d).
 \tag{D.74.9}
\]
This is a discrete continuous $E$-module.  A continuous
map from the compact module $\Delta$ to a finite
discrete module factors through a finite quotient;
the compact-open topology on this Hom space is
discrete, and each such map has an open stabilizer.
In particular ordinary continuous group cohomology
with coefficients $\mathscr H_M$ is defined.

The relative groups used here are, in every positive
degree, the following specified ordinary groups:
\[
 R^n(E,M):=H^n(E,\{D_j\};M)
       =H^{n-1}(E,\mathscr H_M)\qquad(n\geq1).
 \tag{D.74.10}
\]
The corresponding formula for $F$ uses precisely
$\Delta|_F$ and the $L_j$ by (D.74.7).
For clarity, applying continuous Hom to (D.74.6) gives
the exact sequence of discrete $E$-modules
\[
 0\longrightarrow M
 \longrightarrow\operatorname{Hom}_{k,\mathrm{cont}}(P,M)
 \longrightarrow\mathscr H_M\longrightarrow0 .
 \tag{D.74.11}
\]
Exactness follows also directly because (D.74.6)
splits as a sequence of compact $k$-vector spaces.
Shapiro on the three summands gives the actual long
exact sequence
\[
 \cdots\longrightarrow H^n(E,M)
 \xrightarrow{\ r\ }\bigoplus_{j=1}^3 H^n(D_j,M)
 \xrightarrow{\ \delta\ }R^{n+1}(E,M)
 \xrightarrow{\ j\ }H^{n+1}(E,M)\longrightarrow\cdots .
 \tag{D.74.12}
\]
Here $r$ consists of literal subgroup restrictions.
These are the augmentation and Hom constructions of
\cite[Definition~2.1, Propositions~2.3--2.4 and~2.10]
{C-WilkesRelativeAccepted2018}.  In particular the spectral sequence
below will be the ordinary extension spectral sequence
for the actual coefficient module (D.74.9).

\emph{The pairing and its equivariance.}
Let $M^\vee=\operatorname{Hom}_k(M,k)$, with
$(e\ell)(m)=\ell(e^{-1}m)$.  There is the actual
$E$-equivariant coefficient pairing
\[
 M\otimes_k\mathscr H_{M^\vee}\longrightarrow\mathscr H_k,
 \qquad
 m\otimes\varphi\longmapsto
                \bigl[d\longmapsto\varphi(d)(m)\bigr].
 \tag{D.74.13}
\]
Its cup products identify, under (D.74.10), with the
ordinary absolute-relative cup products.  This
identification can be made on the completed bar
resolution: tensor that resolution over $k$ with
$\Delta$, with diagonal group action, and use the
continuous adjunction
\[
 \operatorname{Hom}_F(P_\bullet(F)\widehat\otimes_k\Delta,M)
 \simeq
 \operatorname{Hom}_F(P_\bullet(F),\mathscr H_M).
 \tag{D.74.14}
\]
The augmentation makes the tensor complex a resolution
of $\Delta$; its terms are projective, since the diagonal
action on a free group-module factor can be untwisted
and compact $k$-vector spaces are projective over $k$.
The Alexander--Whitney diagonal followed by (D.74.13)
is the stated cup product.  This also proves compatibility
with the coefficient sequence (D.74.11) and its connecting
maps.  Equivalently these compatibilities are the
natural cup constructions in
\cite[Section~3, especially Propositions~3.1 and~3.5]
{C-WilkesRelativeAccepted2018}.

The full profinite pair $(F,\{L_j\})$ is the completion
of the ordinary pair of pants with its three boundary
subgroups.  Its mod-two duality is the actual cup
duality
\[
 \begin{split}
 H^0(F,M)\otimes H^2(F,\{L_j\};M^\vee)&\longrightarrow k,\\
 H^1(F,M)\otimes H^1(F,\{L_j\};M^\vee)&\longrightarrow k,
 \end{split}
 \tag{D.74.15}
\]
and both pairings are perfect.  One direct comparison
is restriction along the dense inclusions
$F(x,y)\to F$ and $\langle c_j\rangle\to L_j$.
For finite coefficient modules these induce the usual
cohomology isomorphisms for a free group and for an
infinite cyclic group.  The relative long exact
sequences give the relative isomorphisms, with the
same cup products by bar restriction.  Thus the
perfect pairings are those of the compact surface
with its boundary, including its original peripheral
maps.  The completed surface-pair statement and its
cup formulation are also
\cite[Proposition~6.16 and Proposition~5.7]
{C-WilkesRelativeAccepted2018}.

All these pairings are $E/F$-equivariant.  Indeed
(D.74.13) is $E$-equivariant, and the target
\[
 H^1(F,\mathscr H_k)
       =H^2(F,\{L_j\};k)=k
 \tag{D.74.16}
\]
has trivial $E/F$-action.  This last assertion is
literal: the target is one-dimensional over $k$,
so its automorphism group is trivial.  Its
fundamental coclass is the boundary-sum coclass;
the conjugations in (D.74.8) act by inversion on
each peripheral exponent and hence act trivially
on every mod-two boundary character.  No geometric
representative for $\alpha$ is needed for this
equivariance.

Put $A=H^0(F,M)$ and $B=H^1(F,M)$.  They are finite
$k[C_2]$-modules.  Equations (D.74.10) and (D.74.15)
give the following \emph{equivariant} identifications:
\[
 H^0(F,\mathscr H_{M^\vee})=B^\vee,\qquad
 H^1(F,\mathscr H_{M^\vee})=A^\vee .
 \tag{D.74.17}
\]
Every higher row is zero, since $\operatorname{cd}_2(F)=1$.

\emph{The two transgressions have equal rank.}
For a finite $k[C_2]$-module $V$, write
$T=1+\sigma$ and
\[
                 \overline V=\ker T/\operatorname{im}T .
 \tag{D.74.18}
\]
The normalized positive-degree cyclic cochain complex
has value group $V$ in every degree and differential
$T$.  Thus $H^a(C_2,V)=\overline V$ for $a\geq1$.
Multiplication by the standard class
$u\in H^1(C_2,k)$ is the identity under these
identifications.  The pairing
\[
              \overline V\otimes\overline{V^\vee}
                         \longrightarrow k
 \tag{D.74.19}
\]
induced by evaluation is perfect.  To check this,
the annihilator of $\ker T$ is $\operatorname{im}T^\vee$,
and the annihilator of $\operatorname{im}T$ is
$\ker T^\vee$.  The radicals on the kernels are
therefore precisely the images being divided out.
On invariant representatives, the cyclic cup formula
is evaluation itself, since every factor $\sigma^a$
acts trivially on such a representative.

Apply the ordinary Hochschild--Serre spectral sequence
to $F\to E\to C_2$, first with coefficient $M$ and
then with coefficient $\mathscr H_{M^\vee}$.
Its $E_2$ terms in the two cases are
\[
 H^a(C_2,H^b(F,M)),\qquad
 H^a(C_2,H^b(F,\mathscr H_{M^\vee})).
 \tag{D.74.20}
\]
Both have exactly the two possibly nonzero rows
$b=0,1$.  In positive columns their rows are,
respectively,
\[
 \begin{array}{c|cc}
       &b=0&b=1\\ \hline
 \text{absolute}&\overline A&\overline B\\
 \text{shifted relative}&\overline{B^\vee}&\overline{A^\vee}.
 \end{array}
 \tag{D.74.21}
\]
The coefficient pairing (D.74.13), on the same filtered
bar resolutions, supplies a product of these spectral
sequences into that for $\mathscr H_k$.  The
Alexander--Whitney product preserves the quotient
filtration and obeys the cochain Leibniz identity;
therefore its $d_2$ obeys that identity as well.
Its products on the fibre rows are exactly
(D.74.15)--(D.74.17).

Only $d_2$ can be nonzero.  In any sufficiently
positive columns denote the two maps by
\[
 d:\overline B\longrightarrow\overline A,\qquad
 d':\overline{A^\vee}\longrightarrow\overline{B^\vee}.
 \tag{D.74.22}
\]
They do not depend on a further column shift, under
the cyclic identifications, because they commute
with multiplication by $u$.  Take an absolute
row-one element $b$ and a shifted-relative row-one
element $a^\vee$, both in positive columns.
Their product on the fibre row two is zero:
$H^2(F,\mathscr H_k)=0$.  The $d_2$ Leibniz
identity in the target fibre row one is consequently
\[
           \langle d b,a^\vee\rangle+
                    \langle b,d'a^\vee\rangle=0 .
 \tag{D.74.23}
\]
By (D.74.19) these are perfect pairings in the
positive quotient columns.  Hence $d'$ is the
transpose of $d$ for the specified pairings, and
they have a common rank $r$.

There are no later differentials.  It follows that
the two stable total dimensions are
\[
 \begin{split}
 \dim H^n(E,M)
       &=\dim\overline A+\dim\overline B-2r,\\
 \dim H^{n-1}(E,\mathscr H_{M^\vee})
       &=\dim\overline{B^\vee}+
                    \dim\overline{A^\vee}-2r
        =\dim H^n(E,M)
 \end{split}
 \tag{D.74.24}
\]
for all sufficiently large $n$.  The degree shift
in the second line is accounted for by the explicitly
proved one-periodicity (D.74.18) of the positive
quotient columns.  More precisely,
multiplication by the actual inflated $u\in H^1(E,k)$
is an isomorphism in sufficiently high degree on
each group in (D.74.12): on the two-row spectral
sequences it is an isomorphism on both surviving
positive-column rows, and hence on their finite
filtrations.  The same proof applies to each $D_j$,
using its open subgroup $L_j$.

\emph{The rank of the actual boundary restriction.}
Suppose now that $M$ has a specified invariant
perfect bilinear pairing with itself, so that
$M\simeq M^\vee$ as $E$-modules.  This will be the
permutation pairing used below.
For precision, take degree-zero localizations by
the actual class $u$:
\[
 \begin{split}
 \mathcal A_M&=H^\bullet(E,M)[u^{-1}]_0,\\
 \mathcal B_M&=\left(\bigoplus_jH^\bullet(D_j,M)\right)
                                                    [u^{-1}]_0,\\
 \mathcal R_M&=R^\bullet(E,M)[u^{-1}]_0 .
 \end{split}
 \tag{D.74.25}
\]
These are finite vector spaces; each is an actual
sufficiently high cohomology group under the
proved multiplication-by-$u$ maps.  All maps
in (D.74.12) commute with these multiplications.
Localization, or equivalently the direct limit
of that exact sequence under those maps, gives
the exact cyclic sequence
\[
 \mathcal A_M\xrightarrow{\ \bar r\ }\mathcal B_M
       \xrightarrow{\ u^{-1}\bar\delta\ }\mathcal R_M
       \xrightarrow{\ \bar j\ }\mathcal A_M .
 \tag{D.74.26}
\]
For example the middle arrow sends a class
$b/u^n$ to $\delta(b)/u^{n+1}$.
Equation (D.74.24) gives
$\dim\mathcal R_M=\dim\mathcal A_M$.
If the ranks of the three arrows in (D.74.26) are
$a,b,c$, exactness says
\[
 \dim\mathcal A_M=a+c,\qquad
 \dim\mathcal B_M=a+b,\qquad
 \dim\mathcal R_M=b+c .
 \tag{D.74.27}
\]
Consequently $a=b$, and therefore
\[
             \operatorname{rank}\bar r
                         =\tfrac12\dim\mathcal B_M .
 \tag{D.74.28}
\]
This is a rank equality for the actual restriction
map with the original boundary labels.

\emph{Brown restriction and the two occurrences of each section.}
The full finite-coefficient Brown restriction theorem
\cite[Theorem~2 and Section~1.1]
{A-KochloukovaZalesskii2007} applies to $E$:
$F$ is an open subgroup of $2$-cohomological
dimension one, and $E$ has no $C_2\times C_2$.
The two-row sequence (D.74.20) shows that
$H^n(E,M)$ is finite for finite $M$.  Thus the
theorem's dense restriction embedding is onto
its Hausdorff target.  The three full normalizers
are exactly those in (D.74.5).  It follows, for
$n>1$, that
\[
 H^n(E,M)\xrightarrow{\ \sim\ }
                       \bigoplus_{i=0}^2H^n(N_E(S_i),M).
 \tag{D.74.29}
\]
The same theorem applies to each $D_j$, whose
open subgroup $L_j$ has dimension one and whose
two section normalizers are the finite sections
themselves.  Its cohomology is finite by its
two-row cyclic spectral sequence.  Consequently
\[
 H^n(D_j,M)\xrightarrow{\ \sim\ }
           H^n(T_{j,0},M)\oplus H^n(T_{j,1},M),
                    \qquad n>1 ,
 \tag{D.74.30}
\]
where $T_{j,0}=\langle b_j\rangle$ and
$T_{j,1}=\langle c_j^{e_2}b_j\rangle$.

Choose, for each of these six occurrences, an actual
$a_{j,\epsilon}\in F$ with
\[
 T_{j,\epsilon}=a_{j,\epsilon}S_{i(j,\epsilon)}
                                    a_{j,\epsilon}^{-1}.
 \tag{D.74.31}
\]
Such an $F$-conjugator exists by D.34, which identifies
$E$- and $F$-conjugacy for sections.  Its coefficient
transport on cochains is the literal map
\[
 b(t_1,\ldots,t_n)\longmapsto
 a_{j,\epsilon}^{-1}
 b(a_{j,\epsilon}t_1a_{j,\epsilon}^{-1},\ldots,
   a_{j,\epsilon}t_na_{j,\epsilon}^{-1}).
 \tag{D.74.32}
\]
Canonical inner conjugation with this coefficient
transport acts identically on group cohomology.
For example, on the homogeneous bar resolution it
is induced by simultaneous right multiplication
of all vertices by $a_{j,\epsilon}^{-1}$.
The bar prism between that map and the identity
has terms
\[
 \sum_{\ell=0}^{n}
 (g_0,\ldots,g_\ell,
   g_\ell a_{j,\epsilon}^{-1},\ldots,
   g_n a_{j,\epsilon}^{-1});
 \tag{D.74.33}
\]
its boundary identity is the required cochain
homotopy after applying equivariant Hom.

Thus (D.74.29)--(D.74.32) identify $\bar r$ in
(D.74.26) with three duplicated restriction maps.
More explicitly, let
\[
 U_i=H^\bullet(N_E(S_i),M)[u^{-1}]_0,\qquad
 V_i=H^\bullet(S_i,M)[u^{-1}]_0,\qquad
 r_i:U_i\longrightarrow V_i
 \tag{D.74.34}
\]
be induced by the actual inclusion $S_i\subset N_E(S_i)$.
The triangle incidence and the coefficient transports
give
\[
 \mathcal A_M\simeq\bigoplus_iU_i,\qquad
 \mathcal B_M\simeq\bigoplus_i(V_i\oplus V_i),
 \qquad
 \bar r((w_i)_i)=((r_iw_i,r_iw_i))_i .
 \tag{D.74.35}
\]
There is no unidentified automorphism of a coefficient
module in this display; its two copies use exactly
(D.74.32).  In particular
\[
 \operatorname{rank}\bar r=\sum_i\operatorname{rank}r_i,
 \qquad
 \dim\mathcal B_M=2\sum_i\dim V_i .
 \tag{D.74.36}
\]
Together with (D.74.28) these equalities force every
$r_i$ to be surjective.

\emph{A finite permutation coefficient detects every nontrivial commuter.}
Fix $i$ and suppose $1\ne c\in C_i$.
Choose an open normal subgroup $N$ of $E$ with
$N\subset F$ and $c\notin N$.  This is possible
by separating $c$ in a finite continuous quotient
and then intersecting its kernel with $F$.
Write
\[
 J=E/N,\qquad Q=F/N,\qquad \bar S_i=S_iN/N,\qquad
                    M=k[J/\bar S_i].
 \tag{D.74.37}
\]
The quotient still has the specified map to $C_2$,
so $\bar S_i$ has order two,
$Q\cap\bar S_i=1$, and $J=Q\bar S_i$.
The basis pairing
$\langle\delta_a,\delta_b\rangle=1$ for $a=b$
and zero otherwise is a perfect $E$-invariant
self-duality of $M$.  Therefore all of
(D.74.25)--(D.74.36) apply to this exact coefficient
module.

For $n\geq1$, the cyclic calculation (D.74.18)
identifies $H^n(S_i,M)$ with the vector space
on those cosets in $J/\bar S_i$ fixed by $S_i$.
Indeed a two-point orbit gives the regular
$k[C_2]$-module, whose invariant vector is
also its norm image, so it contributes zero.
A singleton orbit contributes its basis vector
in every positive degree.  This identification
commutes with multiplication by $u$ and hence
also describes $V_i$.

The coset $\bar S_i$ is fixed by $S_i$.
The element $\bar c\in Q$ is nonidentity and
centralizes $\bar S_i$, so $\bar c\bar S_i$
is another fixed coset, different from
$\bar S_i$.  Consequently the nonzero class
of $\delta_{\bar S_i}$ in $V_i$ is moved by
$c$, to the different class of
$\delta_{\bar c\bar S_i}$.

On the other hand the image of the restriction
$r_i:U_i\to V_i$ is contained in $C_i$-invariants.
For $c\in C_i$, the cochain transport
$b(g_1,\ldots,g_n)\mapsto
c\,b(c^{-1}g_1c,\ldots,c^{-1}g_nc)$
is inner on $N_E(S_i)$ and acts identically on that
group's cohomology, by (D.74.33) with $a=c^{-1}$.
On $S_i$ its group arguments are unchanged, so it
acts there precisely by $c$ on $M$.
Naturality of restriction
therefore gives the asserted invariance
also after localization.
The class of $\delta_{\bar S_i}$ is not invariant,
contradicting the surjectivity of $r_i$ proved
in (D.74.36).  This proves $C_i=1$.  Equation
(D.74.5) then gives $N_E(S_i)=S_i$, as claimed.

\emph{The exact real-projectivity consequence.}
Every involution in $E$ is conjugate to one of the
$s_i$, so its full centralizer is its own order-two
subgroup.  Also $F$ is an actual open projective
subgroup of index two.  The set of involutions is
closed: it is $\{t\in E:t^2=1\}\cap(E\setminus F)$,
with $E\setminus F$ clopen.  These verify all the
hypotheses of Haran's characterization
\cite[Proposition~2.2]{C-HaranArtinSchreier1993}.
Thus $E$ is real projective.  In the formulation of
\cite[Definition~1.1 and Theorem~2.1]{R-PalQuick2025},
this has the following precise meaning.  For finite
$A,B$, a surjection $\pi:B\to A$, and a continuous
$\varphi:E\to A$, assume that each involution $t$
with $\varphi(t)\ne1$ has an involution $b_t\in B$
with $\pi(b_t)=\varphi(t)$.  There is then a
continuous homomorphism
\[
                 \gamma:E\longrightarrow B,
                         \qquad \pi\gamma=\varphi .
 \tag{D.74.38}
\]
This statement keeps the given finite quotient maps;
it does not require or conclude surjectivity of
$\gamma$.

\emph{Uniqueness of the actual coboundary.}
D.34 supplies $h\in F'$ with
$f=h^{-1}\alpha(h)$.  If $h'$ is any other
solution in $F$, then
\[
 \alpha(h'h^{-1})=h'f(hf)^{-1}=h'h^{-1}.
 \tag{D.74.39}
\]
The element on the right belongs to $C_F(s)=1$,
so $h'=h$.  For $\alpha\in GT_0$, D.44 supplies
a solution with $h\theta(h)=1$.  Uniqueness shows
that this is the same element, with no remaining
fixed-subgroup adjustment.

Keeping the exact original cusps and every unit
exponent, D.35's conjugator set is now parametrized
by the solutions $(\mu,u)$ of
\[
 \begin{gathered}
 \mu\in\widehat{\mathbb Z}^{\times},\qquad u\in F,\\
             x^\mu h y^\mu h^{-1}
                         =u(xy)^\mu u^{-1}.
 \end{gathered}
 \tag{D.74.40}
\]
Its associated maps are the actual maps
$x\mapsto x^\mu$, $y\mapsto hy^\mu h^{-1}$
of D.35.  The closed-subgroup test there proves
that every solution gives a full automorphism;
its converse gives every peripheral conjugator.
The proof above has removed the $C_F(s)$ variable
from that equation.  It has not established the
existence of a pair $(\mu,u)$ satisfying it.
\end{proof}

\paragraph{Scope and proof record.}
The finite coefficient used at the last step is the
actual coset permutation module (D.74.37); no
pro-$2$ quotient is substituted for $F$ or $E$.
The two transgressions in (D.74.22) are not assumed to
vanish.  Their equality of ranks is proved by the
specified cup pairing, and the relative-to-absolute
dimension comparison is then passed through the
actual relative long exact sequence.  The argument
uses D.34's three global section classes and its
labelled triangle, including the fact that each
class occurs twice.  The resulting full self-centralizing property and
real-projectivity are proved above.  Recognition and
the existence in (D.74.40) require their own subsequent maps.
